\chapter{Звёздные достопримечательности}

Отдельного упоминания заслуживают шесть объектов, ради которых
многие туристы едут на~Гран-Канарию.

\section{\sightt{Роке-Нубло} --- «Облачная скала»}

\factline{1813~м. GC-600, парковка 1\,$\frac{1}{2}$~км ниже вершины.}

Вулканическая игла высотой около 80~м, поднимающаяся с~плоского
плато в~центре острова. Для гуанчей --- священное место; сегодня ---
главный символ Гран-Канарии. Восхождение по~маркированной тропе
занимает 45~мин в~одну сторону; обзор с~вершины охватывает
бо́льшую часть острова. См.~маршрут \textbf{GC-1}, стр.~\pageref{hike:gc1}.

\section{\sightt{Дюны Маспаломаса}}

\factline{Охраняемая территория с~1987~г. Вход свободный.}

404~га движущихся дюн, уникальных для Европы. Пески, вопреки
популярному мнению, принесены не~из~Сахары, а~образованы морскими
раковинами и~коралловым песком. Лучшее время --- за~час до~заката,
когда тени подчёркивают рельеф. См.~маршрут \textbf{GC-6}.

\section{\sight{Барранко-де-Гуаядеке}}

Глубокое ущелье на~востоке, где с~доиспанских времён сохранилась
традиция пещерных жилищ. В~деревне \textit{Cueva Bermeja} все дома,
церковь и~даже ресторан вырублены в~туфе. Ущелье --- также ботанический
резерват: здесь цветёт эндемичный \textit{Echium decaisnei guadequianum}.

\section{\sight{Пуэрто-де-Моган}}

Небольшой рыбацкий порт, превращённый в~«игрушечную» курортную
деревню: беленые фасады, цветущие бугенвиллеи, узкие каналы.
Лучшие \textit{papas arrugadas} на~острове --- у~восточного мола.

\section{\sight{Jardín Canario Viera y~Clavijo}}

Крупнейший ботанический сад Испании (27~га) в~долине Гвинигуады,
посвящённый эндемикам Макаронезии. Вход бесплатный. Лучшее время
--- февраль-март, цветение эндемичных \textit{echium}.

\section{\sight{Куэва-Пинтада} --- «Раскрашенная пещера»}

В~городе \place{Гальдар}, на~месте столицы гуанчского королевства.
Сохранились уникальные геометрические росписи доиспанской эпохи;
современный археологический парк закрывает пещеру стеклянным
павильоном.
